%!TEX encoding = UTF8
%!TEX root = 

%%%%%%%%%%%%%%%%%%%%%%%%
%%%%%%%%%% SET 1 %%%%%%%%%%
%%%%%%%%%%%%%%%%%%%%%%%% 


\exe{}{
	Soit $f$ une fonction réelle sur $\R$.
	Montrer que l'ensemble des périodes de $f$ suivant est un groupe commutatif lorsque muni de l'addition $+$ usuelle.
		\[ P_f = \bigset{ p \in \R \tq f(x+p) = f(x) ~ \forall x\in\R } \]
}{exe:groupe-periodes}{
	
}

\exe{}{
	Soit $A$ un anneau unitaire. 
	Montrer que $\bigpar{ A^\times,  \cdot  }$ est un groupe.
}{exe:units-group}{
	todo
}


\exe{difficulty=1}{
	Considérons $ABC$ un triangle équilatéral dans le plan.
	Soit $r$ la rotation centrale qui envoie $A$ sur $B$, $B$ sur $C$, et $C$ sur $A$.
	Soit $s$ la symétrie axiale qui envoie $A$ sur $B$, $B$ sur $A$, et qui laisse $C$ invariant.
	Notons les éléments $r$ et $s$ par leur action sur $A, B, C$ comme suit :
		\begin{align*}
			r = \begin{pmatrix} A&B&C\\B&C&A\end{pmatrix} && \et && s = \begin{pmatrix}A&B&C\\B&A&C\end{pmatrix}.
		\end{align*}
	
	Soit $G = \japb{r, s}$ le groupe généré par $r$ et par $s$ (c'est-à-dire tous les produits de $r$ et de $s$ possibles), doté de la loi de composition naturelle des fontions :
	pour tout $g, h\in G$ et $P$ un point du plan, 
		\[ (gh)(P) = g(h(P)). \]
	\begin{enumerate}
		\item
		Montrer que $G$ a au plus $6$ éléments.
		\item
		Nommer et expliciter tous les éléments du groupe $G$. Ce groupe est appelé \mbox{« groupe diédral »} et noté $D_6$.
		\item
		Décrire la loi de composition de $G$ à l'aide d'une table de multiplication.
	\end{enumerate}
	
	\emph{Remarque : un groupe est entièrement décrit par sa table de multiplication.}
}{exe:D6}{
	todo
}


\exe{difficulty=1}{
	Considérons $\R[X]$, l'ensemble des polynômes à coefficients réels, comme $\R$-espace vectoriel.
	Soit $E = \End(\R[X])$ l'ensemble des endomorphismes (applications $\R$-linéaires) de $\R[X]$ muni de l'opérateur de composition $\circ$.
	\begin{enumerate}
		\item
		Montrer que $E$ est stable pour $\circ$ ; que la loi $\circ$ est associative ; et qu'elle admet un élément neutre.
		\item
		Montrer que l'élément $m_x : p(x) \mapsto x\cdot p(x)$ admet un inverse à gauche mais pas à droite.
		\item
		Montrer que l'élément $\partial : p(x) \mapsto p'(x)$ admet un inverse à droite mais pas à gauche.
		\item
		Donner un élément inversible de $E$ qui n'est pas l'élément neutre.
	\end{enumerate}
	\emph{Conclusion : $E$ n'est pas un groupe.}
}{exe:inv-gauche-not-droite}{
	Il y a en fait une infinité d'inverses à gauche et à droite !
}

\exe{difficulty=1}{
	Soit $G$ un sous-groupe de $\bigpar{\R ,+}$.
	L'ensemble $G$ est dit \emphindex{dense} dans $\R$ toute boule de $\R$ s'intersecte avec $G$. 
	Autrement dit, pour tout $x\in\R$ et tout $\epsilon>0$, il existe un $g\in G$ à distance $\epsilon$ de $x$ : $|x - g| < \epsilon$.
	
	En considérant le plus petit nombre strictement positif de $G$, montrer que $G$ est soit dense, soit de la forme $g\Z$ avec $g\in G$.
	
	\emph{
	Remarque : les nombres représentables par un ordinateur forment un sous-groupe additif de $\bigpar{\R, +}$.
	Chaque nombre représentable est donc multiple entier du plus petit nombre représentable.
	}
}{exe:subgroup-R-p}{
	todo
}


\exe{}{
	Montrer que $u \cdot 0 = 0$ pour tout $u$ élément d'un anneau $R$.
	
	\emph{À gauche : $u$ multiplié par l'élément neutre additif 0.}
	
	\emph{À droite : l'élément neutre additif.}
}{exe:utimes0}{
	Nous avons la suite d'égalités suivantes.
		\begin{align*}
			u \cdot 0 &= u \cdot (0+0)  && \text{car $\bigpar{R, +}$ est un groupe d'élément neutre $0$} \\
				&= u \cdot 0 + u \cdot 0 && \text{par associativité de $ \cdot $ sur $+$} \\
			0 &= u \cdot 0 && \text{en soustrayant $u \cdot 0$}
		\end{align*}
}

\exe{}{
	Montrer que $u \cdot (-1) = -u$ pour tout $u$ élément d'un anneau unitaire $R$.
	
	\emph{À gauche : $u$ multiplié par l'inverse additif de l'élément neutre multiplicatif 1.}
	
	\emph{À droite : l'inverse additif de $u$.}
}{exe:utimes-1}{
	Nous avons la suite d'égalités suivantes.
		\begin{align*}
			u \cdot (-1) + u &= u \cdot (-1+1) && \text{par associativité de $ \cdot $ sur $+$} \\
				&= u \cdot 0 && \\
				&= 0 &&
		\end{align*}
}


\exe{}{
	Considérons $(\Z, +, \cdot)$, l'anneau des entiers relatifs.
	Montrer que $I = 6\Z$ est un idéal de l'anneau. C'est-à-dire :
	\begin{enumerate}
		\item
		$(I, +)$ est un groupe commutatif ; et
		\item
		$n \cdot I \subseteq I$ pour tout $n\in\Z$.
	\end{enumerate}
}{exe:6Z-ideal}{

}

\exe{}{
	Trouver des entiers relatifs $u, v \in \Z$ tels que
		\[ 6u + 15v = 3, \]
	puis montrer que $6\Z + 15\Z = 3\Z$.
}{exe:6Z-p-15Z}{

}

\exe{}{
	Montrer que $6\Z \cap 15\Z = 30\Z$.
}{exe:6Z-cap-15Z}{


}

\exe{}{
	Montrer d'abord que $15\Z \subseteq 5\Z$.
	
	Montrer en suite, plus généralement, que, pour $a, b\in\Z$,
		\begin{align*}
			a | b && \iff && b\Z \subseteq a\Z.
		\end{align*}
	\emph{Conclusion : les relations de division et d'inclusion d'idéaux sont liées mais l'ordre est inversé.}
}{exe:divide-include-Z}{
	Un multiple de 15 est un multiple de 5 donc $15\Z = 5(3\Z) \subseteq 5\Z$.
	
	\underline{Direction $\implies$.}
	
	Si $a|b$ dans $\Z$, $b = aq$ avec $q\in\Z$, et donc $b\Z = a(q\Z) \subseteq a\Z$.
	
	\underline{Direction $\impliedby$.}
	
	Si $b\Z \subseteq a\Z$, alors $b \in a\Z$ et donc $b = aq$ avec un $q\in\Z$.
}


\exe{}{
	Soit $A$ un anneau unitaire.
	Montrer que si $u\in A$ admet un inverse à gauche et un inverse à droite, alors ces inverses sont égaux.
		\[ u^G \cdot u = 1 = u \cdot u^D \implies u^G = u^D. \]
}{exe:inv-gauche-droite}{
	Par associativité, $u^D = (u^G u) u^D = u^G (u u^D) = u^G$.
}

\exe{difficulty=1}{
	Considérons $A = \bigset{ f : \R \rightarrow \R }$ l'ensemble des applications de  $\R$ dans $\R$.
	Nous équipons $A$ d'un addition $+$ et d'une multiplication $\cdot$ définies par
		\begin{align*}
		(f+g)(x) = f(x) + g(x) && \et && (f\cdot g)(x) = f(x)g(x).
		\end{align*}
	\begin{enumerate}
		\item Donner les éléments neutres pour l'addition et la multiplication dans l'anneau $\bigpar{A, +, \cdot}$.
		\item Décrire les diviseurs de zéro de $A$.
		\item Montrer que chaque élément non nul est soit un diviseur de zéro, soit inversible.
	\end{enumerate}
}{exe:ring-of-maps}{
	todo
}

\exe{}{
	Soit $A$ un anneau unitaire.
	Montrer qu'une unité $u\in A^\times$ n'est pas un diviseur de zéro.
}{exe:unit-not-zero-divisor}{
}

\exe{difficulty=0}{
	Considérons le sous-ensemble $E$ des matrices réelles $2\times2$ suivant.
		\[ E = \left\{ \begin{pmatrix} a & -b \\ b & a \end{pmatrix} \tq a, b \in \R \right\}. \]
	Montrer que $E$ est un corps lorsque muni de l'addition et de la multiplication usuelles des matrices.
}{exe:C-as-matrix}{
	todo
}

\exe{difficulty=1}{
	Soit $p\in\N$ un nombre premier.
	Montrer que $\F_p = \bigpar{ \Z/p\Z, +, \cdot }$ est un corps.
	
	Trouver l'inverse de 6 dans $\F_{15}$ à l'aide de l'exercice \ref{exe:6Z-p-15Z}.
}{exe:Fp-field}{
	todo
}

\exe{difficulty=1}{
	Soit $(A, +, \cdot)$ un anneau commutatif unitaire et $\circ$ la loi définie par
		\[ a \circ b = a + b - a\cdot b. \]
	\begin{enumerate}
		\item 
		Montrer que $\circ$ est associative d'élément neutre 0.
		\item
		Montrer que $\bigpar{A-\{1\}, \circ}$ est un groupe si et seulement si $A$ est un corps.
	\end{enumerate}
}{exe:circle-op}{
	todo
}


\exe{}{
	Montrer que tout corps est intègre : aucun élément non nul ne divise zéro.
}{exe:field-integral}{

}


%%%%%%%%%%%%%%%%%%%%%%%%
%%%%%%%%%% SET 2 %%%%%%%%%%
%%%%%%%%%%%%%%%%%%%%%%%% 

\exe{}{
	Considérons $\Z, \R[X], \C[X], \Q[X]$ les anneaux avec addition et multiplication usuelles.
	\begin{enumerate}
		\item
		Énumérer les unités de $\Z$.
		\item
		Décrire les irréductibles de $\Z$.
		\item
		Décrire les inversibles de $\R[X]$.
		\item
		Montrer que $X^2 + 1$ est irréductible dans $\R[X]$.
		\item
		Montrer que $6X^2 + 6$ est irréductible dans $\R[X]$.
		\item
		Montrer que $X^2 + 1$ est réductible dans $\C[X]$.
		\item
		Donner un polynôme de $\Q[X]$ qui est irréductible dans $\Q[X]$ mais réductible dans $\R[X]$
	\end{enumerate}
}{exe:irr-in-Z-Rx-Cx}{
	\begin{enumerate}
		\item
		Les unités sont les éléments inversibles multiplicativement : $\Z^\times = \{ +1 ; -1 \}$.
		\item
		Les $\pm p$ sont irréductibles car si $\pm p = a \times b$, alors $a$ ou $b$ est unité.
		Les autres entiers relatifs sont réductibles.
		\item
		Les inversibles de $\R[X]$ sont les constantes non nulles : si $pq = 1$ dans $\R[X]$, alors $p$ et $q$ sont nécessairement constants et non nuls car $\deg(pq) = \deg(p) + \deg(q) = 0$.
		\item
		Si $X^2 + 1 = pq$ avec $p, q$ des diviseurs propres (donc non constants), alors $p$ et $q$ sont linéaires et admettent une racines dans $\R$.
		Ce n'est bien sûr pas le cas de $X^2 + 1$.
		\item
		Si $6X^2 + 6$ se réduit, alors $6X^2 + 6 = pq$, et $X^2 + 1 = (\frac16 p)q$.
		Comme $X^2 + 1$ est irréductible, ce n'est pas possible.
		\item
		$X^2 + 1 = (X+i)(X-i)$ dans $\C[X]$.
		\item
		$X^2 - 2 = (X+\sqrt2)(X-\sqrt2)$ fonctionne.
	\end{enumerate}
}

\exe{}{
	Soit $A$ un anneau commutatif unitaire et $u\in A^\times$ un élément inversible de $A$.
	Montrer que
	\begin{align*}
		\text{$a$ est irréductible} && \iff && \text{$u\cdot a$ est irréductible}.
	\end{align*}
}{exe:irr-up-to-unit}{
	\underline{Direction $\implies$.}
	
	Si $a$ est irréductible, il n'a aucun facteur propre (ni unité, ni $a$ fois une unité).
	Supposons par contradiction que $u\cdot a$ soit réductible :
		\[ u \cdot a = p \cdot q, \]
	avec $p, q \in A$ deux facteurs propres.
	
	Alors, comme $u$ est inversible,
		\[ a = (u^{-1} p) \cdot q, \]
	ce qui contredit l'irréducibilité de $a$.
	
	\underline{Direction $\impliedby$.}
	
	Si $u \cdot a$ est irréductible, alors $u^{-1} (u \cdot a)$ l'est par la direction $\implies$ déjà traitée.
	Donc $a$ est irréductible.
}

\exe{}{
	Montrer que l'anneau $(\Z, +, \cdot)$ est principal : tout idéal $I$ est généré par un seul élément.
	
	\emph{Indication : prendre un $n\in I$ de valeur absolue minimale non nul et montrer que $I = n\Z$.}
}{exe:Z-is-principal}{
	Soit $I\subseteq\Z$ un idéal.
	Prenons un élément $n \in I-\{0 \}$ de valeur absolue minimale (il y en a deux : $n$ et $-n$).
	Montrons que $I = \japb{n} = n\Z$.
	
	\underline{Inclusion $\supseteq$.}
	
	Il est clair que $n\Z \subseteq I$ car $(I,+)$ est un groupe : l'élément neutre $0$ appartient à $I$, et comme $n\in I$, la stabilité de $+$ et l'existence d'un inverse additif implique que $n+n = 2n \in I$, et $-n \in I$. 
	En répétant les sommes, $n\Z \subseteq I$.
	
	\underline{Inclusion $\subseteq$.}
	
	Soit $x\in I$. La division euclidienne de $x$ par $n$ donne
		\[ x = nq + r, \]
	avec $|r| < |n|$ (la valeur absolue est ajoutée car nous somme dans $\Z$ et que $x$ peut être un nombre négatif !).
	Comme $I$ est un idéal, $nq \in I$ par stabilité par multiplication externe, et $x - nq \in I$ par stabilité de l'addition interne.
	Par suite, $r\in I$.
	Comme $n$ a été choisi de valeur absolue minimale non nul, nous avons nécessairement $r=0$ et donc $x \in n\Z$.
}


\exe{}{
	Soit $A$ un anneau commutatif et $a\in A$ générant un idéal $I$ :
		\[ I = \japb{a}. \]
	On rappelle que $I$ est le plus petit idéal contenant $a$.
	
	Montrer que
		\[ I = a \cdot A  = \bigset{ a \cdot x \tq x \in A }. \]
	\emph{Exemple : dans $\Z$, $\japb{15} = 15\Z$.}
}{exe:gen-as-sum-of-ideals}{
	Soit $J$ un idéal contenant $a$.
	Montrons que $J$ contient $a\cdot A$ et que $a \cdot A$ est un idéal : ceci montrera que $\japb{a} = a\cdot A$.
	
	D'une part, par stabilité externe d'un idéal, $a\in J$ implique $a \cdot A \subseteq J$.
	
	D'autre pat, $a\cdot A$ est un idéal car $(a \cdot A, +)$ est un groupe (élément neutre, stabilité, et inverse sont les seules propriétés à vérifier) et car $a\cdot A$ est stable par multiplication par un élément de $A$ : $a \cdot A \cdot x \subseteq a\cdot A$ pour tout $x\in A$.
}

\exe{}{
	Soit $A$ un anneau commutatif et certains éléments $a_1, a_2, \dots, a_n \in A$ générant un idéal $I$ :
		\[ I = \japb{a_1, a_2, \dots, a_n}. \]
	On rappelle que $I$ est le plus petit idéal contenant $a_1, a_2, \dots, a_n$.
	
	Montrer que
		\[ I = a_1 \cdot A + a_2 \cdot A + \dots + a_n \cdot A. \]
	\emph{Exemple : dans $\Z$, $\japb{78, 66, 24} = 78\Z + 66\Z + 24\Z$.}
}{exe:gen-as-sum-of-ideals2}{
	Analogue à \ref{exe:gen-as-sum-of-ideals} : tout idéal contenant $a_1, \dots, a_n$ contient nécessairement les multiples de ces éléments ainsi que les sommes.
	Il suffit alors de montrer que $a_1 A + \dots + a_n A$ est un idéal, ce qui n'est pas trop difficile.
}

\exe{}{
	Soit $A$ un anneau commutatif et certains éléments $a_1, a_2, \dots, a_n \in A$.
	Montrer que pour tout idéal $I$,
	\begin{align*}
		\japb{ a_ 1, a_2, \dots, a_n } \subseteq I && \iff && a_1 \in I, a_2 \in I, \dots, a_n \in I.
	\end{align*}
	\emph{Conclusion : un idéal est inclus dans un autre dès que tous ses générateurs le sont.}
}{exe:gen-incl-ideal-incl}{
	La direction $\implies$ est claire car les générateurs appartiennent à l'idéal généré par ceux-ci.
	
	Pour la direction $\impliedby$, supposons $a_1 \in I, a_2 \in I, \dots, a_n \in I$.
	Alors, comme $I$ est stable par multiplication externe par un élément de $A$, nous avons
		\[ a_1 \cdot A \subseteq I, a_2 \cdot A \subseteq I, \cdots, a_n \cdot A \subseteq I. \]
	De plus, comme $I$ est stable par addition interne, nous avons
		\[ a_1 A + a_2 A + \dots + a_n A \subseteq I, \]
	comme requis.
}

\exe{}{
	Soit $A$ un anneau commutatif unitaire et $u\in A^\times$ un élément inversible de $A$.
	Montrer que, pour tout $a\in A$,
		\[ \japb{ a } = \japb{u \cdot a} \]
	\emph{ Conclusion : multiplier un générateur par un inversible ne change pas l'idéal généré.}
}{exe:ideal-up-to-unit}{
	D'après l'exercice \ref{exe:gen-incl-ideal-incl}, il suffit de montrer que le générateur du premier idéal est inclus dans le deuxième, et inversement.
	Or, clairement, $u\cdot a \in \japb{a}$, et $a \in \japb{u \cdot a}$ car $a = u^{-1} (u\cdot a)$.
}

\exe{}{
	Soit $A$ un anneau commutatif et $a, b\in A$ deux éléments.
	Montrer que
	\begin{align*}
		a|b && \iff && \japb{b} \subseteq \japb{a}.
	\end{align*}
}{exe:divide-include}{
	\underline{Direction $\implies$.}
	
	Si $a|b$, alors $b=a\cdot q$ avec $q\in A$. Ainsi $b \in \japb{a} = a \cdot A$ donc $\japb{b} \subseteq \japb{a}$ par inclusion du générateur (voir exercice \ref{exe:gen-incl-ideal-incl}.)
	
	\underline{Direction $\impliedby$.}
	
	Si $\japb{b}\subseteq \japb{a}$, alors $b \in \japb{b}$ et donc $b \in \japb{a} = a \cdot A$.
	Il existe donc un $q\in A$ tel que $b = aq$.
}


\exe{}{
	Montrer l'équivalence suivante dans $\Z$.
	\begin{align*}
		a = bq + r && \iff && \begin{pmatrix} 0 & 1 \\ 1 & -q \end{pmatrix} \begin{pmatrix} a \\ b \end{pmatrix} = \begin{pmatrix} b \\ r \end{pmatrix}.
	\end{align*}
}{exe:div-euclide-matrix}{
	Calcul matriciel.
}

\exe{}{
	Calculer $\pgcd(78, 66)$ à l'aide de l'algorithme d'Euclide.
	Utiliser la formulation matricielle pour trouver $u, v\in\Z$ tels que
		\[ 78u + 66v = \pgcd(78, 66). \]
}{exe:Bezout-Z}{
	Les divisions euclidiennes successives sont les suivantes.
	\begin{align*}
		78 &= 66\times 1 + 12 \\
		66 &= 12 \times 5 + 6 \\
		12 &= 6 \times 2 + 0
	\end{align*}
	Le dernier reste non nul est 6, c'est le PGCD.
	
	En formulation matricielle, l'exercice \ref{exe:div-euclide-matrix} donne
	\begin{align*}
		\begin{pmatrix} 0 & 1 \\ 1 & -2 \end{pmatrix} \begin{pmatrix} 0 & 1 \\ 1 & -5 \end{pmatrix} \begin{pmatrix} 0 & 1 \\ 1 & -1 \end{pmatrix}  \begin{pmatrix} 78 \\ 66 \end{pmatrix} &=  \begin{pmatrix} 6 \\ 0 \end{pmatrix}, \\
		\begin{pmatrix} -5 & 6 \\ 11 & -13 \end{pmatrix} \begin{pmatrix} 78 \\ 66 \end{pmatrix} &=  \begin{pmatrix} 6 \\ 0 \end{pmatrix}.
	\end{align*}
	La première ligne donne :
		\[ 78\times(-5) + 66\times(6) = 6, \]
	qu'on vérifiera bien sûr !
}

\exe{}{
	Soit $A$ un anneau commutatif unitaire principal, et $a, b\in A$ deux éléments.
	Nous savons donc que $\japb{a, b} = \japb{d}$ avec $d\in A$.
	
	Montrer que $d$ vérifie
	\begin{enumerate}
		\item
		$d | a$ et $d|b$ ; et
		\item
		si $d'|a$ et $d'|b$, alors $d|d'$.
	\end{enumerate}
	\emph{Conclusion : le PGCD existe et est unique à multiplication par unité près dans un anneau principal.}
}{exe:gcd-principal-ideal}{
	\begin{enumerate}
		\item
		Comme $\japb{a, b} = \japb{d}$, alors $a \in \japb{d} = d\cdot A$ et donc $d$ divise $a$.
		De façon identique, $d$ divise $b$.
		\item
		Si $d'|a$ et $d'|b$, alors $d'\in aA + bA = \japb{a, b} = \japb{d} = d\cdot A$.
		Il suit donc que $d$ divise $d'$.
	\end{enumerate}
}

\exe{}{
	Trouver $u, v\in\Z$ tels que $18u + 23v = 1$.
	
	En déduire la valeur de $18^{-1}$ dans $\F_{23}$ (le corps des entiers modulo 23).
}{exe:Bezout-18-23}{
	Les divisions successives
	\begin{align*}
		23 &= 18 \times 1 + 5, \\
		18 &= 5 \times3 + 3, \\
		5 &= 3 \times 1 + 2, \\
		3 &= 2 \times 1 + 1 \\
		2 &= 1 \times 2 + 0
	\end{align*}
	Donne la relation matricielle
	\begin{align*}
		\begin{pmatrix} 0 & 1 \\ 1 & -2 \end{pmatrix}\begin{pmatrix} 0 & 1 \\ 1 & -1 \end{pmatrix}\begin{pmatrix} 0 & 1 \\ 1 & -1 \end{pmatrix}\begin{pmatrix} 0 & 1 \\ 1 & -3 \end{pmatrix}\begin{pmatrix} 0 & 1 \\ 1 & -1 \end{pmatrix}
		 \begin{pmatrix} 23 \\ 18 \end{pmatrix} &=  \begin{pmatrix} 1 \\ 0 \end{pmatrix}, \\
		\begin{pmatrix} -7 & 9 \\ 18 & -23 \end{pmatrix} \begin{pmatrix} 23 \\ 18 \end{pmatrix} &=  \begin{pmatrix} 1 \\ 0 \end{pmatrix}.
	\end{align*}
	La première ligne donne donc $23 \times (-7) + 18 \times (9) = 1$.
	
	Modulo $23$, cette relation donne $18 \times 9 = 1 \mod 23$ et donc l'inverse de 18 est 9 dans $\F_{23}$.
}

\exe{}{
	Poser la division euclidienne de
	\begin{enumerate}[itemsep=10pt]
		\item
		$X^4 + 2X + 1$ par $2X^2 - 1$ dans $\Q[X]$ ;
		\item
		$X^3 + 5X^2 + 2X$ par $\pi X^2 + 3$ dans $\R[X]$ ;
		\item
		$X^2 + 7X + 2$ par $18X + 1$ dans $\F_{23}[X]$ ;
	\end{enumerate}
}{exe:euclidean-div-Qx-Cx-Fp}{
	La division euclidienne fonctionne en annulant le terme dominant du dividende en lui soustrayant des multiples du diviseur.
	Pour $a(x) = a_n X^n + \dots$ et $b(X) = b_m X^m + \dots$ et $m < n$, alors l'opération
		\[ a - b \times (a_n b_m^{-1} X^{n-m}) \]
	permet d'annuler le coefficient de $a$, avant de répéter.
	\begin{enumerate}
		\item
		En posant, 
		\polylongdiv[style=D]{X^4 + 2X + 1}{2X^2 - 1}
		
		donc $(X^4 + 2X + 1) = (2X^2 - 1) (\frac12 X^2 + \frac14) + (2X + \frac54)$.
		\item
		D'abord,
			\[ X^3 + 5X^2 + 2X - \frac1{\pi}X\left(\pi X^2 + 3\right) = 5X^2 + \left(2-\frac3{\pi}\right)X, \]
		puis
			\[  5X^2 + \left(2-\frac3{\pi}\right)X - \frac5{\pi} \left(\pi X^2 + 3\right) = \left(2-\frac3{\pi}\right)X  - \frac{15}{\pi}. \]
		et donc finalement
			\[ X^3 + 5X^2 + 2X = \left(\frac1{\pi} X + \frac5{\pi}\right) \left(\pi X^2 + 3\right) + \left(2-\frac3{\pi}\right)X  - \frac{15}{\pi}. \]
		\item 
		L'inverse de $18$ est $9$ dans $\F_{23}$ d'après l'exercice \ref{exe:Bezout-18-23}, donc 
			\[ X^2 + 7X + 2 - (18X+1)(9X) = -2X + 2 \mod 23\]
		puis 
			\[ -2X + 2 + 2(18X+1)(9) = 20 \mod 23 \]
		et donc finalement
			\[ X^2 + 7X + 2 = (18X+1)(9X - 18) + 20 \mod 23. \]
	\end{enumerate}
}


\exe{difficulty=1}{
	Soit $p\in\N$ un premier impair.
	\begin{enumerate}
		\item
		Montrer que le polynôme $ax^2 + bx + c \in \Z[x]$ avec $p\nmid a$ est réductible dans $\F_p[x]$ si et seulement si le discriminant $\Delta = b^2 - 4ac$ est un carré parfait modulo $p$.
		\item 
		En déduire un polynôme de $\Z[x]$ irréductible dans $\F_5[x]$ mais réductible dans $\F_7[x]$.
		\item
		Montrer que $x^2 + x + 1$ est irréductible dans $\F_2[x]$.
	\end{enumerate}
}{exe:discr-Fp}{
	Le polynôme est de degré 2 et est donc réductible si et seulement s'il admet une racine dans $\F_p$.
	Les équations suivantes sont équivalentes et admettent une solution si et seulement si $\Delta$ est un carré parfait.
	\begin{align*}
		ax^2 + bx + c = 0 \\
		4a^2x^2 + 4abx + 4ac = 0 \\
		(2ax+b)^2 - b^2 + 4ac = 0 \\
		(2ax+b)^2 = \Delta
	\end{align*}
	Nous avons utilisé ici que $p\neq2$ pour avoir que $4$ est inversible dans $\F_p$, et que $p\nmid a$ pour que $a$ soit inversible dans $\F_p$.

	Les carrés parfaits de $\F_5$ sont $\bigset{ 0 ; 1 ; 4 }$.
	Ceux de $\F_7$ sont $\bigset{ 0 ; 1 ; 4 ; 2 }$.
	Obtenir $\Delta = 2$ est donc notre unique chance et le polynôme $x^2 - 2$ un candidat évident.
	Celui-ci est irréductible dans $\F_5[x]$ mais se décompose comme $(x-3)(x+3)$ dans $\F_7[x]$.
	
	Il existe des tables de polynômes irréductibles pour chaque corps fini, dans lesquelles nous pouvons trouver 
	\begin{itemize}[label=$\bullet)$]
		\item $x^2 + x + 1$, de discriminant 2 dans $\F_5$ (irréductible) et 3 dans $\F_7$ (irréductible) ;
		\item $x^2 + x + 2$, de discriminant 3 dans $\F_5$ (irréductible) et 0 dans $\F_7$ (réductible) ;
		\item $x^2 + 4x + 1$, de discriminant 2 dans $\F_5$ (irréductible) et 5 dans $\F_7$ (irréductible) ;
		\item etc...
	\end{itemize}
}


\exe{}{
	Montrer que le polynôme $x^2 - 1$ se décompose de deux façons différentes dans $\sfrac{\Z}{8\Z}[x]$.
	
	\emph{Conclusion : l'anneau $\sfrac{\Z}{8\Z}[x]$ n'est pas factoriel car la décomposition n'est pas unique, même à permutation ou multiplication par constante près.}
}{exe:Zo8Zx-non-factoriel}{
	L'élément 1 admet quatre racines carrés dans $\Z/8\Z$ !
	Celles-ci sont $\pm1$ et $\pm3$.
	Ainsi $x^2 - 1 = (x+1)(x-1) = (x+3)(x-3)$ dans l'anneau.
}


\exe{difficulty=1}{
	Soient $a, b \in\Z[i]$, $b\neq0$.
	Montrer qu'il existe $q, r\in\Z[i]$ vérifiant $a = bq+r$ et $|r| < |b|$.
}{exe:Zi-euclidien}{
	Le ratio $\frac{a}{b}$ appartient à $\Q[i] = \{p + qi \tq p, q \in \Q \}$.
	Notons donc $\frac{a}b = p + qi$ avec $p, q \in \Q$, et $u, v$ les entiers les plus proches de $p, q$ respectivement.
	Autrement dit, $|p-u| \leq \frac12$ et $|q-v| \leq \frac12$.
	
	Avec ce choix, 
		\[ \left| \frac{a}b  - (u+vi) \right|^2 = (p-u)^2 + (q-v)^2 \leq \frac14 + \frac14 = \frac12. \]
	En multipliant par $|b|$, on trouve
		\[ | a - (u+vi) b | \leq \frac12 |b|, \]
	ce qui conclut en posant $q=u+vi$ et $r = a-qb$ vérifiant $|r| \leq \frac12|b| < |b|$.
}



\exe{}{
	Le polynôme $f(x) = x^4+x^3+x^2+x+1$ est-il irréductible dans $\F_5[x]$ ?
}{exe:irreductible-F5}{
	$f(1) = 0$ dans $\F_5$ donc $x-1$ divise $f(x)$ dans $\F_5[x]$.
	En posant la division euclidienne, on trouve
		\[ x^4 + x^3+ x^2 + x + 1 = (x-1)(x^3 + 2x^2 + 3x + 4) \mod 5. \]
}



